%!TEX root = ../Thesis.tex
\section{Einleitung}
\subsection{Problemstellung}

Der Einsatz von Chatbots in Unternehmen hat in den vergangenen Jahren deutlich an Bedeutung gewonnen. Insbesondere durch die Fortschritte im Bereich leistungsfähiger Large Language Models (LLMs) haben sich Chatbots zu einem zentralen Bestandteil der digitalen Unternehmensinfrastruktur entwickelt. Sie können sowohl Kunden als auch Mitarbeitende in unterschiedlichen Anwendungsszenarien unterstützen, indem sie Informationen bereitstellen, Prozesse automatisieren oder Handlungsempfehlungen ableiten.

Auch das Unternehmen Phoenix Contact hat im Zuge dieser Entwicklungen einen Enterprise-Chatbot in den operativen Betrieb integriert. Dieser ist bereits in der Lage, Nutzeranfragen strukturiert zu verarbeiten und dabei Informationen aus internen Wissensbasen zu berücksichtigen. Trotz dieser Funktionalitäten zeigt sich in der Praxis, dass die Antworten des Chatbots überwiegend standardisiert erfolgen. Individuelle Nutzerpräferenzen werden dabei nicht sitzungsübergreifend oder über längere Konversationen hinweg von dem Chatbot beachtet.

Ein wesentlicher Grund hierfür liegt darin, dass vergangene Interaktionen nicht nachhaltig verarbeitet werden. Präferenzen und Anforderungen, die Nutzer in früheren Gesprächen äußern, werden weder systematisch extrahiert noch langfristig gespeichert oder bei späteren Konversationen berücksichtigt. Dadurch wird jede Sitzung isoliert betrachtet, sodass kein langfristiges Verständnis des jeweiligen Nutzers aufgebaut werden kann. 

Dieses Problem führt dazu, dass die Erwartungen der Nutzer an die Leistungsfähigkeit des Systems nicht vollständig erfüllt werden können. Gerade aus der Sicht des Kunden kann ein ineffizientes Chatbot-System negative Einflüsse auf die Geschäftsbeziehung oder die Wahrnehmung des Unternehmens haben. Gleichzeitig beeinträchtigt es die interne Produktivität der Mitarbeiter, da Präferenzen redundant geäußert werden müssen, wodurch der Aufwand in einzelnen Konversationen steigt. 

Vor diesem Hintergrund wird deutlich, dass die langfristige Berücksichtigung von Nutzerpräferenzen sowohl aus Kunden- als auch aus Mitarbeiterperspektive von zentraler Bedeutung ist. Um dies zu gewährleisten, müssen geeignete Mechanismen entwickelt werden, die Nutzerpräferenzen automatisiert erfassen, strukturiert speichern und in zukünftige Interaktionen integrieren können.

\subsection{Zielsetzung}
Daher befasst sich diese Arbeit mit der Personalisierung eines Chatbots im Unternehmenskontext. Ziel der Arbeit ist die Entwicklung eines Konzepts, das es ermöglicht, Nutzeranfragen unter Berücksichtigung individueller Präferenzen personalisiert zu verarbeiten. Dabei liegt der Fokus insbesondere auf der systematischen Erfassung, Speicherung und Nutzung von Nutzerpräferenzen, die im Rahmen von Konversationen entstehen.
Zur Strukturierung der Untersuchung werden im Rahmen dieser Arbeit folgende Forschungsfragen betrachtet:

\begin{itemize}
	\item Wie können Nutzerpräferenzen aus Konversationen systematisch extrahiert werden?
	\item Wie können Nutzerpräferenzen in zukünftigen Konversationen integriert werden?
	\item Kann eine systematische Integration von Nutzerpräferenzen die Antwortgenerierung und die Nutzerzufriedenheit nachweisbar beeinflussen?
\end{itemize}



\subsection{Vorgehensweise}

Im ersten Teil der Arbeit werden die theoretischen Grundlagen dargestellt. In diesem Teil der Arbeit soll vor allem ein grundlegendes Verständnis für Chatbots, Agenten, Large Language Models und Vektordatenbanken geschaffen werden, sodass Design Entscheidungen hinsichtlich der Zielsetzung nachvollziehbar sind. Gleichzeitig wird genauer erläutert welche Arten von Nutzerpräferenzen im Rahmen dieser Arbeit berücksichtigt werden. Darauf aufbauend wird der aktuelle Stand der Forschung untersucht. Dabei liegt der Fokus insbesondere auf bestehenden Ansätzen, die für die Personalisierung von Chatbots existieren.

Im anschließenden Praxisteil wird der Bezug zur betrieblichen Problemstellung bei Phoenix Contact hergestellt. Dazu wird zunächst der aktuelle Entwicklungsstand des eingesetzten Chatbots analysiert, um bestehende Prozessabläufe sowie mögliche Schwachstellen zu identifizieren

